%% kapa: exam=1 round=23 year=2012 period=2 form=A subject=appraisal_law subject_ko=부동산관계법규 seq=30 number=70
%% source: Q-Net 감정평가사 기출문제 (cst00309) artlSeq=5052580 file=raw/1008/5052580/2012년 제23회 감정평가사 제1차시험 기출문제(2교시 A형).pdf
%% conversion: kordoc 4.24.51 -> tools/gen1.py (auto); review=proofread-sonnet; flags=none
\begin{question}{70}
건축법령상 건축허가 및 신고에 관한 설명으로 옳지 않은 것은?
\begin{choices}
  \item 허가권자는 주변의 교육환경을 고려할 때 부적합하다고 인정되는 숙박시설에 대해서는 건축위원회의 심의를 거쳐 건축허가를 거부할 수 있다.
  \item 허가권자는 공사에 착수하였으나 공사의 완료가 불가능하다고 인정되는 경우 건축허가를 취소할 수 있다.
  \item 연면적 200\,m$^2$ 미만이고 3층 미만인 건축물의 대수선은 건축신고 사항이다.
  \item 건축신고일부터 1년 이내에 공사에 착수하지 아니하면 그 신고의 효력은 없어진다.
  \item 국토해양부장관은 국토관리를 위하여 특히 필요하다고 인정할 경우 건축허가를 받은 건축물의 착공을 제한할 수 있다.
\end{choices}
\end{question}
\answer{2}
